% EDITING: type inside the final {...} of each field. Change the shown font size
% (for example, font=10pt) in that same field block when needed.

\GrantSection{5}{Human Subjects, Ethics, Inclusion, and Regulatory Strategy}

\GrantGuidance{Use this section as a drafting companion to the required human-subjects and clinical-trials forms. It does not replace an IRB protocol, informed-consent document, FDA submission, single-IRB plan, or agency-specific study record.}

\begin{tabularx}{\linewidth}{X X X}
\textbf{Human subjects involved?} & \textbf{NIH-defined clinical trial?} & \textbf{Delayed onset?}\\[-1pt]
\GrantTableField[id=hs-human-subjects,font=11pt,width=0.95\linewidth,maxchars=24]{Yes} &
\GrantTableField[id=hs-nih-trial,font=11pt,width=0.95\linewidth,maxchars=24]{Yes} &
\GrantTableField[id=hs-delayed-onset,font=11pt,width=0.95\linewidth,maxchars=24]{No}\\[6pt]
\textbf{IND / IDE status} & \textbf{IRB / sIRB status} & \textbf{ClinicalTrials.gov status}\\[-1pt]
\GrantTableField[id=hs-ind-ide,font=7.5pt,width=0.95\linewidth,maxchars=80]{IND prepared; IDE/SR-NSR pending} &
\GrantTableField[id=hs-irb,font=7.5pt,width=0.95\linewidth,maxchars=80]{IRB in preparation; sIRB N/A} &
\GrantTableField[id=hs-ctgov,font=7.5pt,width=0.95\linewidth,maxchars=80]{PRS package setup pending}
\end{tabularx}

\GrantTextArea[
  id=hs-risks-protections, font=10pt, height=2.95in, maxchars=1800
]{Risk to participants and adequacy of protections}{Participants face more than minimal risk because the intervention combines a first-in-human perioperative targeted agent strategy, major pancreatic surgery, anesthesia, intensive perioperative monitoring, and a governed robotic/AI operating environment \citep{phase1ind,onpremwhippl}. Foreseeable risks include drug toxicity, surgical bleeding, leak, infection, delayed gastric emptying, thromboembolic events, device malfunction, software or sensor failure, workflow interruption, privacy ^^Jloss, and burdens associated with additional study procedures. Protections are layered: eligibility limits enrollment to adults with resectable or borderline-resectable KRAS G12-mutated PDAC and adequate organ function; all patient-facing actions occur under licensed surgeon and anesthesiology authority; predefined verification, validation, and uncertainty-quantification gates must clear before robot-enabled execution; immediate human override and emergency-stop authority is continuously available; pharmacy, operating-room, and intensive-care SOPs are prespecified; adverse events are rapidly graded, attributed, and reviewed; enrollment pauses are triggered for treatment-related death, unexpected grade 4 or 5 toxicity, major uncontrolled cyber-physical failure, or sponsor/IRB/medical-monitor concern \citep{cfr312paiuni,cfr050paiuni,ichpaistduni,clintrustpai,vvuqllmverif,h2pancrevvuq}. ^^JAdditional protections include credentialed site activation, independent medical monitoring, source-data review, protocol deviation tracking, and withdrawal rights without penalty.
}
\GrantTextArea[
  id=hs-benefits-knowledge, font=10pt, height=2in, maxchars=1400
]{Potential benefits, importance of knowledge, and risk-benefit assessment}{Direct clinical benefit cannot be promised, and some participants may receive no therapeutic benefit beyond standard perioperative management. Potential direct benefits include tumor-control or operative advantages sufficient to support successful resection, improved perioperative coordination, and closer protocol-driven monitoring; however, these are hypotheses rather than established outcomes in this first-in-human setting \citep{phase1ind,onpremwhippl}. The knowledge value is substantial because the study addresses whether a targeted therapy, robotic procedure, and verification-first AI governance layer can be introduced into oncology research under accountable regulatory, ethical, and operational controls. The risk-benefit assessment is therefore acceptable only with strict eligibility, staged activation, intensive oversight, rapid stopping rules, and transparent governance. In a disease with poor prognosis and limited curative opportunity, the proposed bounded Phase 1 design is justified as a safety and feasibility study, not as a claim of superior efficacy \citep{daraxwhipple,clintrustpai,nih_rfa_rm_27_001}.
}
\GrantTextArea[
  id=hs-consent, font=10pt, height=2.45in, maxchars=1650
]{Informed consent / assent, capacity, and additional safeguards for vulnerable populations}{Informed consent will be obtained by qualified study personnel before any research-only procedure, using an IRB-approved consent form written in plain language and a process that gives prospective participants adequate time to review the study, ask questions, discuss alternatives, and decide free of coercion \citep{cfr050paiuni}. Because the initial target population is competent adults undergoing high ^^Jrisk cancer surgery, assent is not expected to apply. Decisional capacity will be assessed clinically, and consent will not be obtained during acute sedation, severe pain crises, or other conditions that materially impair understanding. Key disclosures will include the first-in-human nature of the program, uncertain benefit, major surgical and drug risks, robotic and software governance features, alternatives to participation, privacy protections, ClinicalTrials.gov posting, and the right to withdraw. Additional safeguards include optional re-consent after material protocol change, witness procedures when needed, translated materials and interpreter support for participants with limited English proficiency, and exclusion of prisoners, children, and other populations requiring additional subpart-specific protections unless a future amendment supplies the required regulatory justifications and safeguards.
}
\GrantTextArea[
  id=hs-inclusion, font=10pt, height=2in, maxchars=1750
]{Inclusion by sex, race, ethnicity, and age; scientific justification for any exclusions}{Enrollment will be open to eligible adults regardless of sex, gender, race, or ethnicity, and recruitment materials, screening logs, and accrual reviews will be used to monitor whether study participation reflects the catchment area served by the participating site \citep{nih_rfa_rm_27_001}. No race- or ethnicity-based exclusions are scientifically justified. The protocol should aim for broad inclusion of women and men and for ^^Jequitable access among historically underrepresented populations through bilingual outreach, navigator support, flexible scheduling, and travel/logistics assistance when feasible. Age inclusion will begin at 18 years and older, with no arbitrary upper-age exclusion; older adults may participate if surgical candidacy and organ function criteria are met. Exclusion of children is scientifically and ethically justified because pancreaticoduodenectomy with perioperative daraxonrasib in KRAS G12-mutated PDAC is an adult first-in-human oncology intervention with disease epidemiology, operative risk, and dose/procedure uncertainty that do not support pediatric enrollment at this stage.
}
\GrantTextArea[
  id=hs-privacy, font=10pt, height=2.1in, maxchars=1450
]{Privacy, confidentiality, HIPAA, Certificate of Confidentiality, and controlled-access strategy}{Participant privacy and confidentiality will be protected through coded study identifiers, separation of direct identifiers from analytic datasets, minimum-necessary data access, role-based permissions, audit-trailed systems, and secure transfer and storage procedures. Research use and disclosure of protected health information will occur under IRB-approved HIPAA authorization or other legally permissible mechanism, with clear notice of what information will be used, who may receive it, and how participants may revoke authorization where applicable. If NIH funding is awarded and the project collects or uses identifiable, sensitive information, the study will be automatically deemed issued a Certificate of Confidentiality; non-NIH-funded linked studies would require separate CoC review if applicable \citep{nih_rfa_rm_27_001}. Only de-identified or suitably transformed datasets will be shared publicly; datasets with residual re-identification risk, operative video, detailed device logs, or linkable biospecimen/genomic variables will use controlled-access release, documented data-use conditions, and repository or enclave controls.
}
\GrantTextArea[
  id=hs-regulatory-pathway, font=10pt, height=2.1in, maxchars=1450
]{Regulatory pathway, required submissions, and responsible parties}{The pathway is a sponsor-investigator drug-plus-device clinical investigation. Before any participant facing activity, the study will submit an IND under 21 CFR Part 312 with Form FDA 1571, protocol, drug information, safety-monitoring materials, and ClinicalTrials.gov certification \citep{phase1ind,cfr312paiuni}. In parallel, the robotic/device component requires a documented 21 CFR Part 812 strategy: FDA-approved IDE if significant risk, or an IRB-supported nonsignificant-risk pathway with abbreviated IDE controls if justified \citep{onpremwhippl}. Parallel submissions include local IRB review, consent and HIPAA documents, PRS registration, monitoring materials, SOPs for software change control and emergency override, protocol amendments, IND safety reporting, continuing review, and final closeout. The sponsor-investigator leads FDA correspondence, monitoring, and overall conduct; the site PI/surgeon of record leads participant care; the IRB provides ethics review; and designated PRS, data, pharmacy, quality, and security staff maintain compliance and records.
}
